% Learning InitechBase, Chapter 10: Keeping Your Data Safe (AN-16).  Every
% panel replayed (examples/l10-*.dot); sources cmd/mutate.c (writes at once,
% m_confirm_file for ZAP, RECALL), cmd/set.c (SAFETY).
\chapter{Keeping Your Data Safe}
\Gates{test-pubs-examples, test-interp-replace, test-samir-safety-emu}
InitechBase writes every change to the disk at once: there is no separate
step that saves your work, and no command that takes a change back.  This
lesson gathers the habits that keep the records of Initech safe while you work
with them.

\begin{ChapterCovers}
\item Working copies
\item Looking before you change
\item Keeping a value before you change it
\item Deleting without losing
\item The confirmation of ZAP
\item Leaving InitechBase properly
\end{ChapterCovers}

\ATEHead{Preparing for This Chapter}
Start InitechBase with a fresh copy of the sample files.

\ATEHead{Working Copies}
Keep the original diskettes in their package and work on copies.  Before a
change to many records, and at the end of each day's work, copy your databases,
their memo files and their index files to a diskette with the InitechDOS
command \file{COPY}.  A database is several files: \file{STAFF.DBF} needs
\file{STAFF.DBT} and its index files with it.

\ATEHead{Look Before You Change}
\cmd{REPLACE}, \cmd{DELETE} and \cmd{RECALL} with \kw{FOR} change every record
that meets the condition.  List the records with the same condition first, and
change them only when the list is what you expect:
\begin{Steps}
\item Type \type{USE STAFF}.
\item Type \type{LIST LASTNAME, SALARY FOR SALARY > 100000}.
\item Type \type{REPLACE SALARY WITH SALARY - 10000 FOR SALARY > 100000}.
\item Type \type{LIST LASTNAME, SALARY FOR LASTNAME = 'PORTER'}.
\end{Steps}
\DotExample{examples/l10-check}

\ATEHead{Keep the Old Value}
When you change one field of one record, store the old value in a memory
variable first.  If the change was wrong, \cmd{REPLACE} puts the old value
back:
\DotExample{examples/l10-undo}

\ATEHead{Deleting Without Losing}
\cmd{DELETE} only marks records.  While the records are marked, \cmd{RECALL}
brings them all back; only \cmd{PACK} removes them.
\DotExample{examples/l10-recall}
Wait until you are sure before you \cmd{PACK}, and copy the database first.

\ATEHead{The Confirmation of ZAP}
\cmd{ZAP} removes every record at once.  While SET SAFETY is ON, it names the
database and asks for confirmation; any answer but Y leaves the records where
they are:
\begin{Steps}[resume]
\item Type \type{USE INVOICE}, then \type{? RECCOUNT()}.
\item Type \type{ZAP} and answer \type{N}.
\item Type \type{? RECCOUNT()}.
\end{Steps}
\DotExample{examples/l10-zapno}
Leave SET SAFETY ON.

\ATEHead{Leaving Properly}
Leave InitechBase with \cmd{QUIT} before you turn off the computer or take out
a diskette.  \cmd{QUIT} closes every database with its index files, so that
the files on the disk agree with one another.

\ATEHead{Where Do You Go From Here?}
You have finished \doctitle{Learning InitechBase}.  \doctitle{Using
InitechBase}, the next section of this volume, describes every command and
function in full, and \doctitle{Programming with InitechBase}, in Volume Two,
teaches you to write programs.

\ATEHead{Summary}
\begin{enumerate}
\item Changes are written to the disk at once; work on copies.
\item List with the same \kw{FOR} condition before you change records.
\item Keep an old value in a memory variable before you replace it.
\item \cmd{RECALL} brings back marked records until \cmd{PACK} removes them.
\item Leave SET SAFETY ON; leave InitechBase with \cmd{QUIT}.
\end{enumerate}
